\subsection{Matrix-Vector Multiplication}
\hrulefill
\subsubsection*{Vector-Matrix Multiplication}
\paragraph*{Definition}
Let $A$ be an $m \times n$ matrix written $A = [a_1 \ a_2 \ \ldots \ a_n]$ where $a_i$ is the $i$-th column of $A$. 
Let $\vec{x}$ be an $n \times 1$ column vector written as $\vec{x} = \begin{bmatrix}
    x_1 \\
    \vdots \\
    x_n
\end{bmatrix}$. Then the matrix-vector product $A\vec{x}$ is defined as
\begin{align*}
    A\vec{x} &= [a_1 \ a_2 \ \ldots \ a_n] \begin{bmatrix}
    x_1 \\
    \vdots \\
    x_n
\end{bmatrix} \\
&= x_1a_1 + x_2a_2 + \ldots + x_na_n \\
&= \sum_{i=1}^n x_ia_i
\end{align*}

\paragraph*{Notes}
\begin{itemize}
    \item If $A$ is an $m \times n$ matrix, then $A\vec{x}$ is defined only for $\vec{x} \in \mathbb{R}^n$.
    \item The inner sizes of the vector and matrix must match for the product to be defined.
    \item $A\vec{x}$ is a linear combination of the columns of $A$ with coefficients $x_i$ in order.
\end{itemize}

\paragraph*{Example} Let $B = \begin{bmatrix}
    1 & 0 & 3 \\
    2 & -1 & 1
\end{bmatrix}$ and $\vec{q} = \begin{bmatrix} 1 \\ 2 \\ 3 \end{bmatrix}$. Find $B\vec{q}$.
\begin{align*}
B\vec{q} &= 1 \cdot \begin{bmatrix} 1 \\ 2 \end{bmatrix} + 2 \cdot \begin{bmatrix} 0 \\ -1 \end{bmatrix} + 3 \cdot \begin{bmatrix} 3 \\ 1 \end{bmatrix} = \begin{bmatrix} 1 \\ 2 \end{bmatrix} + \begin{bmatrix} 0 \\ -2 \end{bmatrix} + \begin{bmatrix} 9 \\ 3 \end{bmatrix} = \begin{bmatrix} 10 \\ 3 \end{bmatrix}
\end{align*}

\paragraph*{Example} Let $A = \begin{bmatrix}
    0 & -1 & 2 \\
    5 & 3 & -4 \\
    1 & 0 & 1
\end{bmatrix}$ and $\vec{p} = \begin{bmatrix} 1 \\ 0 \\ -3 \end{bmatrix}$. Find $A\vec{p}$.
\begin{align*}
A\vec{p} &= 1 \cdot \begin{bmatrix} 0 \\ 5 \\ 1 \end{bmatrix} + 0 \cdot \begin{bmatrix} -1 \\ 3 \\ 0 \end{bmatrix} + (-3) \cdot \begin{bmatrix} 2 \\ -4 \\ 1 \end{bmatrix} = \begin{bmatrix} 0 \\ 5 \\ 1 \end{bmatrix} + \begin{bmatrix} 0 \\ 0 \\ 0 \end{bmatrix} + \begin{bmatrix} -6 \\ 12 \\ -3 \end{bmatrix} = \begin{bmatrix} -6 \\ 17 \\ -2 \end{bmatrix}
\end{align*}

\paragraph*{Theorem}
A system of linear equations $\begin{cases}
a_{11}x_1 + a_{12}x_2 + \ldots + a_{1n}x_n = b_1 \\
a_{21}x_1 + a_{22}x_2 + \ldots + a_{2n}x_n = b_2 \\
\vdots \\
a_{m1}x_1 + a_{m2}x_2 + \ldots + a_{mn}x_n = b_m
\end{cases}$ can be written in matrix form as $A\vec{x} = \vec{b}$ where
\[A = \begin{bmatrix}
    a_{11} & a_{12} & \ldots & a_{1n} \\
    a_{21} & a_{22} & \ldots & a_{2n} \\
    \vdots & \vdots & \ddots & \vdots \\
    a_{m1} & a_{m2} & \ldots & a_{mn}
\end{bmatrix}, \quad \vec{x} = \begin{bmatrix}
    x_1 \\
    x_2 \\
    \vdots \\
    x_n
\end{bmatrix}, \quad \vec{b} = \begin{bmatrix}
    b_1 \\
    b_2 \\
    \vdots \\
    b_m
\end{bmatrix}\] 
$A$ is called the coefficient matrix, $\vec{x}$ is called the variable vector, and $\vec{b}$ is called the constant vector. The system of equations has a solution if and only if there exists a vector $\vec{x}$ such that $A\vec{x} = \vec{b}$.
\begin{itemize}
    \item Every system of linear equations can be written as $A\vec{x} = \vec{b}$ 
    \item $A\vec{x} = \vec{b}$ is consistant if and only if $\vec{b}$ is a linear combination of columns of $A$.
\end{itemize}

\subsubsection*{Identitiy Matrix}
\paragraph*{Definition}
The $n \times n$ identity matrix, denoted by $I_n$, is the square matrix with 1's on the main diagonal and 0's elsewhere. That is,
\[I_n = \begin{bmatrix}
    1 & 0 & \ldots & 0 \\
    0 & 1 & \ldots & 0 \\
    \vdots & \vdots & \ddots & \vdots \\
    0 & 0 & \ldots & 1
\end{bmatrix}\]

\paragraph*{Properties of the Identity Matrix}
\begin{itemize}
    \item rank($I_n$) = $n$
    \item $I_n\vec{x} = \vec{x}$ for all $\vec{x} \in \mathbb{R}^n$
\end{itemize}

\paragraph*{More Indentities}
Let $A$ and $B$ be $n \times n$ matricies and $\vec{x}, \vec{y} \in \mathbb{R}^n$
\begin{itemize}
    \item $A(\vec{x} + \vec{y}) = A\vec{x} + A\vec{y}$
    \item $A(k\vec{x}) = k(A\vec{x})$ for $k \in \mathbb{R}$
    \item $(A + B)\vec{x} = A\vec{x} + B\vec{x}$
\end{itemize}

\paragraph*{Theorem}
Let $A\vec{x} = \vec{b}$ be a system of linear equations and rank($A$) = $r$. Then the 
rank of the augmented matrix $[A | \vec{b}]$ is either $r$ or $r + 1$. 
The system is consistent if and only if rank($A$) = rank($[A | \vec{b}]$).

\paragraph*{Example} Let $A = \begin{bmatrix}
    1 & 0 & 2 & 0 \\
    0 & 1 & 3 & 0 \\
    0 & 0 & 0 & 0 \\
    0 & 0 & 0 & 0
\end{bmatrix}$ and $\vec{b} = \begin{bmatrix} 1 \\ 2 \\ 1 \\ 0 \end{bmatrix}$. Determine if the system $A\vec{x} = \vec{b}$ is consistent.
\begin{align*}
    \text{rank}(A) = 2 \text{ and rank}([A | \vec{b}]) = 3 \implies \text{the system is inconsistent}
\end{align*}

\paragraph*{Dot Product}
Let $\vec{a} = \begin{bmatrix}
    a_1 \\
    \vdots \\
    a_n
\end{bmatrix}$ and $\vec{b} = \begin{bmatrix}
    b_1 \\
    \vdots \\
    b_n
\end{bmatrix}$ be two vectors in $\mathbb{R}^n$. The dot product of $\vec{a}$ and $\vec{b}$, denoted by $\vec{a} \cdot \vec{b}$, is defined as
\[\vec{a} \cdot \vec{b} = a_1b_1 + \ldots + a_nb_n = \sum_{i=1}^n a_ib_i\]

The result of the dot product is a scalar.

\paragraph*{Definition}
Let $A$ be an $m \times n$ matrix and $\vec{x} \in \mathbb{R}^n$, then the $i$th entry of the vector $A\vec{x}$ is 
the dot product of the $i$th row of $A$ with $\vec{x}$

\paragraph*{Example} 2Let $A = \begin{bmatrix}
    1 & -1 & 3 & -2 \\
    2 & 3 & 5 & 7 \\
    3 & 0 & 4 & 0 \\
    4 & 5 & 2 & 0
\end{bmatrix}$ and $\vec{x} \begin{bmatrix} 0 \\ 1 \\ -1 \\ 5 \end{bmatrix}$. Find the second entry of $A\vec{x}$ without computing the entire product.
\[
P_2 = [2 \ 3 \ 5 \ 7] \cdot \begin{bmatrix} 0 \\ 1 \\ -1 \\ 5 \end{bmatrix} = 2(0) + 3(1) + 5(-1) + 7(5) = 3 - 5 + 35 = 33
\]

\paragraph*{Example} Let $C = \begin{bmatrix}
    2 & -1 & 3 & 5 \\
    0 & 2 & -3 & 1  \\
    -3 & 4 & 1 & 2 
\end{bmatrix}$ and $\vec{q} = \begin{bmatrix} 2 \\ 1 \\ 0 \\ -1 \end{bmatrix}$. Find $C\vec{q}$.
\begin{align*}
C\vec{q} &= 2 \cdot \begin{bmatrix} 2 \\ 0 \\ -3 \end{bmatrix} + 1 \cdot \begin{bmatrix} -1 \\ 2 \\ 4 \end{bmatrix} + 0 \cdot \begin{bmatrix} 3 \\ -3 \\ 1 \end{bmatrix} + (-1) \cdot \begin{bmatrix} 5 \\ 1 \\ 2 \end{bmatrix} \\
&= \begin{bmatrix} 4 \\ 0 \\ -6 \end{bmatrix} + \begin{bmatrix} -1 \\ 2 \\ 4 \end{bmatrix} + \begin{bmatrix} 0 \\ 0 \\ 0 \end{bmatrix} + \begin{bmatrix} -5 \\ -1 \\ -2 \end{bmatrix} \\
&= \begin{bmatrix} -2 \\ 1 \\ -4 \end{bmatrix}
\end{align*}

\paragraph*{Standard Basis Vectors}
The standard basis vectors in $\mathbb{R}^n$ are the vectors $e_1, e_2, \ldots, e_n$ where $e_i$ has a 1 in the $i$-th position and 0's elsewhere. That is,
\[e_1 = \begin{bmatrix}
    1 \\
    0 \\
    \vdots \\
    0
\end{bmatrix}, \quad e_2 = \begin{bmatrix}
    0 \\
    1 \\
    \vdots \\
    0
\end{bmatrix}, \quad e_i = \begin{bmatrix}
    0 \\
    \vdots \\
    0 \\
    1 \\
    0 \\
    \vdots \\
    0
\end{bmatrix}\]

We can use these basis vectors to isolate the columns of a matrix. For example, if $A = [a_1 \ a_2 \ \ldots \ a_n]$ is an $m \times n$ matrix, then
\begin{align*}
Ae_1 &= a_1 \\
Ae_2 &= a_2 \\
&\vdots \\
Ae_n &= a_n
\end{align*} 

\paragraph*{Theorem}
$A, B \in \mathbb{R}^{m \times n}, \vec{x} \in \mathbb{R}^n \implies A=B \implies \forall i < n \in \mathbb{N}, Ae_i = Be_i$

\paragraph*{Note}
$A\vec{x} = B\vec{x}$ does not imply $A = B$ for all $\vec{x} \in \mathbb{R}^n$. For example, if $A$ and $B$ are both $m \times n$ matrices with all entries equal to 0, then $A\vec{x} = B\vec{x}$ for all $\vec{x} \in \mathbb{R}^n$, but $A$ and $B$ are not necessarily the same matrix.

\paragraph*{Theorem}
$A\vec{x} = B\vec{x} \implies \forall \vec{x} \in \mathbb{R}^n, A\vec{x} = B\vec{x}$

\hrulefill

\subsubsection*{Transformations}
\paragraph*{Definition}
Functions $T : \mathbb{R}^n \to \mathbb{R}^m$ are called transformations. 
The set $\mathbb{R}^n$ is called the domain of $T$ and the set $\mathbb{R}^m$ is called the codomain of $T$. 
The image of a vector $\vec{x} \in \mathbb{R}^n$ under the transformation $T$ is denoted by $T(\vec{x})$. 

\paragraph*{Matrix Transformations}
A matrix transformation is denoted by $T_A(\vec{x}) = A\vec{x}$ for some matrix $A_{m \times n}$. If $A \in \mathbb{R}^{m \times n}$, then $T_A : \mathbb{R}^n \to \mathbb{R}^m$

\paragraph*{Example} $\forall \vec{x} \in \mathbb{R}^2, T_A(\vec{x}) = A\vec{x}$ and let $A = \begin{bmatrix}
    1 & -3 \\
    3 & 5 \\
    -1 & 7
\end{bmatrix}$ and $\vec{u} = \begin{bmatrix} 2 \\ -1 \end{bmatrix}$. Find $T_A(\vec{u})$.
\[
T_A(\vec{u}) = \begin{bmatrix}
    1 & -3 \\
    3 & 5 \\
    -1 & 7
\end{bmatrix} \begin{bmatrix} 2 \\ -1 \end{bmatrix} = \begin{bmatrix} 1(2) + (-3)(-1) \\ 3(2) + 5(-1) \\ -1(2) + 7(-1) \end{bmatrix} = \begin{bmatrix} 5 \\ 1 \\ -9 \end{bmatrix}
\]

\paragraph*{Example}
Let $\vec{b} = \begin{bmatrix}
    3 \\ 2 \\ -5
\end{bmatrix}$ and $\vec{c} = \begin{bmatrix}
    3 \\ 2 \\ 5
\end{bmatrix}$. Find $\vec{x} \in \mathbb{R}^2 \ni T_A(\vec{x}) = \vec{b}$ and $\vec{y} \in \mathbb{R}^2 \ni T_A(\vec{y}) = \vec{c}$.
\begin{align*}
    A\vec{x} = \vec{b} &\implies \begin{bmatrix}
        1 & -3 \\
        3 & 5 \\
        -1 & 7
    \end{bmatrix} \begin{bmatrix}
        x_1 \\
        x_2
    \end{bmatrix} = \begin{bmatrix}
        3 \\
        2 \\
        5
    \end{bmatrix} \implies \begin{bmatrix}
        1 & -3 & | & 3 \\
        3 & 5 & | & 2 \\
        -1 & 7 & | & -5
    \end{bmatrix} \xmapsto{RREF} \begin{bmatrix}
        1 & 0 & | & \frac{3}{2} \\
        0 & 1 & | & -\frac{1}{2} \\
        0 & 0 & | & 0
    \end{bmatrix} \\
    \implies & \vec{x} = \begin{bmatrix}
        \frac{3}{2} \\
        -\frac{1}{2}
    \end{bmatrix} \\
    \\
    A\vec{y} = \vec{c} &\implies \begin{bmatrix}
        1 & -3 | 3 \\
        3 & 5 | 2 \\
        -1 & 7 | 5
    \end{bmatrix} \xmapsto{RREF} \begin{bmatrix}
        1 & -3 & | & 3 \\
        0 & 1 & | & -\frac{1}{2} \\
        0 & 0 & | & 10
    \end{bmatrix} \implies \text{the system is inconsistent}
\end{align*}